
\documentclass{article}
\usepackage{colortbl}
\usepackage{makecell}
\usepackage{multirow}
\usepackage{supertabular}

\begin{document}

\newcounter{utterance}

\centering \large Interaction Transcript for game `hot\_air\_balloon', experiment `air\_balloon\_survival\_mt\_complexity\_hard', episode 3 with gemma{-}4{-}26b{-}a4b{-}it.
\vspace{24pt}

{ \footnotesize  \setcounter{utterance}{1}
\setlength{\tabcolsep}{0pt}
\begin{supertabular}{c@{$\;$}|p{.15\linewidth}@{}p{.15\linewidth}p{.15\linewidth}p{.15\linewidth}p{.15\linewidth}p{.15\linewidth}}
    \# & $\;$A & \multicolumn{4}{c}{Game Master} & $\;\:$B\\
    \hline

    \theutterance \stepcounter{utterance}  
    & & \multicolumn{4}{p{0.6\linewidth}}{
        \cellcolor[rgb]{0.9,0.9,0.9}{
            \makecell[{{p{\linewidth}}}]{
                \texttt{\tiny{[P1$\langle$GM]}}
                \texttt{Int qiegħed tipparteċipa f'logħba ta' negozjar kollaborattiva.} \\
\\ 
\texttt{Flimkien ma' parteċipant ieħor, trid taqblu fuq sett wieħed ta' oġġetti li se jinżammu. Kull wieħed minnkom għandu l{-}fehma tiegħu dwar kemm kull oġġett hu importanti għalih (importanza). Int ma tafx kif il{-}parteċipant l{-}ieħor jivvaluta l{-}oġġetti. Barra minn hekk, tingħata l{-}isforz li jitlob kull oġġett.} \\
\texttt{Tistgħu taqblu fuq sett biss jekk l{-}isforz totali tal{-}oġġetti magħżula ma jaqbiżx limitu kondiviż:} \\
\\ 
\texttt{LIMITU = 5638} \\
\\ 
\texttt{Hawn huma l{-}valuri individwali tal{-}isforz tal{-}oġġetti:} \\
\\ 
\texttt{Isforz tal{-}oġġett = \{"C08": 110, "C55": 727, "A43": 527, "A67": 855, "B42": 942, "A62": 617, "B19": 301, "C43": 130, "C73": 916, "C10": 212, "B36": 146, "A18": 559, "C61": 933, "C67": 740, "C42": 33, "C70": 799, "B40": 324, "C39": 841, "A21": 925, "C16": 639\}} \\
\\ 
\texttt{Hawn hi l{-}fehma personali tiegħek dwar l{-}importanza ta' kull oġġett:} \\
\\ 
\texttt{Valuri tal{-}importanza tal{-}oġġetti = \{"C08": 204, "C55": 821, "A43": 621, "A67": 949, "B42": 1036, "A62": 711, "B19": 395, "C43": 224, "C73": 1010, "C10": 306, "B36": 240, "A18": 653, "C61": 1027, "C67": 834, "C42": 127, "C70": 893, "B40": 418, "C39": 935, "A21": 1019, "C16": 733\}} \\
\\ 
\texttt{Għan:} \\
\\ 
\texttt{L{-}għan tiegħek hu li tinnegozja sett kondiviż ta' oġġetti li jibbenefikak kemm jista' jkun (jiġifieri, jimmassimizza l{-}importanza totali GĦALIK), filwaqt li tibqa' fil{-}LIMITU. M'intix obbligat tagħmel PROPOSTA f'kull messaġġ {-} tista' sempliċement tinnegozja wkoll. It{-}tattiċi kollha huma permessi!} \\
\\ 
\texttt{Protokoll ta' Interazzjoni:} \\
\\ 
\texttt{Tista' tuża biss il{-}formati strutturati li ġejjin f'messaġġ:} \\
\\ 
\texttt{PROPOSTA: \{'A', 'B', 'C', …\}} \\
\texttt{Ipproponi li jinżammu eżattament dawk l{-}oġġetti.} \\
\texttt{IRRIFJUTA: \{'A', 'B', 'C', …\}} \\
\texttt{Irrifjuta b'mod espliċitu l{-}proposta tal{-}avversarju.} \\
\texttt{ARGUMENT: \{'...'\}} \\
\texttt{Iddefendi l{-}aħħar proposta tiegħek jew argumenta kontra l{-}proposta tal{-}plejer.} \\
\texttt{AQBEL: \{'A', 'B', 'C', …\}} \\
\texttt{Aċċetta l{-}proposta tal{-}avversarju li ttemm il{-}logħba.} \\
\\ 
\\ 
\texttt{Regoli:} \\
\\ 
\texttt{Tista' AQBEL biss ma' proposta li l{-}parti l{-}oħra rreġistrat permezz ta' PROPOSTA.} \\
\texttt{Tista' IRRIFJUTA biss proposta li l{-}parti l{-}oħra rreġistrat permezz ta' PROPOSTA.} \\
\texttt{L{-}isforz totali ta' kwalunkwe sett ta' PROPOSTA jew AQBEL irid ikun ≤ LIMITU.} \\
\texttt{TIŻVELAX il{-}punteġġi moħbija tal{-}importanza tiegħek.} \\
\texttt{Tikketta f'format strutturat trid tkun segwita minn kolon u spazju vojt. L{-}argument irid ikun sett ta' python li fih 0 jew aktar strings.} \\
\texttt{Għalhekk, irid ikun fil{-}forma "TIKKETTA: \{...\}"} \\
\texttt{Segwi b'mod strett il{-}protokoll ta' interazzjoni u TIKTIBX xejn lil hinn mill{-}istruttura mogħtija.} \\
\texttt{Il{-}logħba tispiċċa meta naħa waħda tagħti AQBEL għal PROPOSTA magħmula mill{-}plejer l{-}ieħor.} \\
\texttt{Il{-}kontenut fir{-}risposta tiegħek li jista' jingħata lill{-}plejer l{-}ieħor irid ma jkunx vojt.} \\
\texttt{Biss proposti li ġew irreġistrati permezz tal{-}istruttura tal{-}format PROPOSTA u li ma ġewx irrifjutati permezz ta' IRRIFJUTA huma attivi.} \\
\\ 
            }
        }
    }
    & \\ \\

    \theutterance \stepcounter{utterance}  
    & \multicolumn{4}{p{0.6\linewidth}}{
        \cellcolor[rgb]{0.9,0.9,0.9}{
            \makecell[{{p{\linewidth}}}]{
                \texttt{\tiny{[P1$\rangle$GM]}}
                \texttt{PROPOSTA: \{'C08', 'C55', 'A43', 'A67', 'B42', 'A62', 'B19', 'C43', 'C10', 'B36', 'C42', 'B40'\}} \\
            }
        }
    }
    & & \\ \\

    \theutterance \stepcounter{utterance}  
    & & \multicolumn{4}{p{0.6\linewidth}}{
        \cellcolor[rgb]{0.9,0.9,0.9}{
            \makecell[{{p{\linewidth}}}]{
                \texttt{\tiny{[GM$\rangle$P2]}}
                \texttt{Qed tipparteċipa f'logħba ta' negozjar kollaborattiva.} \\
\\ 
\texttt{Flimkien ma' parteċipant ieħor, trid taqblu fuq sett wieħed ta' oġġetti li se jinżammu. Kull wieħed minnkom għandu l{-}fehma tiegħu dwar kemm kull oġġett hu importanti għalih (importanza). Ma tafx kif il{-}parteċipant l{-}ieħor jivvaluta l{-}oġġetti. Barra minn hekk, tingħata l{-}isforz li jitlob kull oġġett.} \\
\texttt{Tistgħu taqblu fuq sett biss jekk l{-}isforz totali tal{-}oġġetti magħżula ma jaqbiżx limitu kondiviż:} \\
\\ 
\texttt{LIMITU = 5638} \\
\\ 
\texttt{Hawn huma l{-}valuri individwali tal{-}isforz tal{-}oġġetti:} \\
\\ 
\texttt{Isforz tal{-}oġġett = \{"C08": 110, "C55": 727, "A43": 527, "A67": 855, "B42": 942, "A62": 617, "B19": 301, "C43": 130, "C73": 916, "C10": 212, "B36": 146, "A18": 559, "C61": 933, "C67": 740, "C42": 33, "C70": 799, "B40": 324, "C39": 841, "A21": 925, "C16": 639\}} \\
\\ 
\texttt{Hawn hi l{-}fehma personali tiegħek dwar l{-}importanza ta' kull oġġett:} \\
\\ 
\texttt{Valuri tal{-}importanza tal{-}oġġetti = \{"C08": 204, "C55": 821, "A43": 621, "A67": 949, "B42": 1036, "A62": 711, "B19": 395, "C43": 224, "C73": 1010, "C10": 306, "B36": 240, "A18": 653, "C61": 1027, "C67": 834, "C42": 127, "C70": 893, "B40": 418, "C39": 935, "A21": 1019, "C16": 733\}} \\
\\ 
\texttt{Għan:} \\
\\ 
\texttt{L{-}għan tiegħek hu li tinnegozja sett kondiviż ta' oġġetti li jibbenefikak kemm jista' jkun (jiġifieri, jimmassimizza l{-}importanza totali GĦALIK), filwaqt li tibqa' fil{-}LIMITU. M'intix obbligat tagħmel PROPOSTA f'kull messaġġ {-} tista' sempliċement tinnegozja wkoll. It{-}tattiċi kollha huma permessi!} \\
\\ 
\texttt{Protokoll ta' Interazzjoni:} \\
\\ 
\texttt{Tista' tuża biss il{-}formati strutturati li ġejjin f'messaġġ:} \\
\\ 
\texttt{PROPOSTA: \{'A', 'B', 'C', …\}} \\
\texttt{Ipproponi li jinżammu eżatt dawk l{-}oġġetti.} \\
\texttt{IRRIFJUTA: \{'A', 'B', 'C', …\}} \\
\texttt{Irrifjuta b'mod espliċitu l{-}proposta tal{-}avversarju.} \\
\texttt{ARGUMENT: \{'...'\}} \\
\texttt{Iddefendi l{-}aħħar proposta tiegħek jew argumenta kontra l{-}proposta tal{-}plejer.} \\
\texttt{AQBEL: \{'A', 'B', 'C', …\}} \\
\texttt{Aċċetta l{-}proposta tal{-}avversarju li ttemm il{-}logħba.} \\
\\ 
\\ 
\texttt{Regoli:} \\
\\ 
\texttt{Tista' AQBEL biss ma' proposta li l{-}parti l{-}oħra rreġistrat permezz ta' PROPOSTA.} \\
\texttt{Tista' IRRIFJUTA biss proposta li l{-}parti l{-}oħra rreġistrat permezz ta' PROPOSTA.} \\
\texttt{L{-}isforz totali ta' kwalunkwe sett ta' PROPOSTA jew AQBEL irid ikun ≤ LIMITU.} \\
\texttt{TIŻVELAX il{-}punteġġi moħbija tal{-}importanza tiegħek.} \\
\texttt{Tikketta f'format strutturat trid tkun segwita b'kolon u spazju. L{-}argument irid ikun sett ta' Python li jkun fih 0 jew aktar strings.} \\
\texttt{Għalhekk, irid ikun fil{-}forma "TIKKETTA: \{...\}"} \\
\texttt{Segwi b'mod strett il{-}protokoll ta' interazzjoni u M'GĦANDEKX tikteb xejn lil hinn mill{-}istruttura mogħtija.} \\
\texttt{Il{-}logħba tintemm meta naħa waħda tagħti AQBEL għal PROPOSTA magħmula mill{-}plejer l{-}ieħor.} \\
\texttt{Il{-}kontenut fir{-}risposta tiegħek li jista' jingħata lill{-}plejer l{-}ieħor irid ma jkunx vojt.} \\
\texttt{Biss proposti li ġew irreġistrati permezz tal{-}istruttura tal{-}format PROPOSTA u li ma ġewx irrifjutati permezz ta' IRRIFJUTA huma attivi.} \\
\\ 
\\ 
\\ 
\texttt{Issa se tirċievi l{-}ewwel messaġġ tal{-}plejer l{-}ieħor.} \\
\\ 
\texttt{PROPOSTA: \{'C08', 'C55', 'A43', 'A67', 'B42', 'A62', 'B19', 'C43', 'C10', 'B36', 'C42', 'B40'\}} \\
            }
        }
    }
    & \\ \\

    \theutterance \stepcounter{utterance}  
    & & & \multicolumn{4}{p{0.6\linewidth}}{
        \cellcolor[rgb]{0.9,0.9,0.9}{
            \makecell[{{p{\linewidth}}}]{
                \texttt{\tiny{[GM$\langle$P2]}}
                \texttt{IRRIFJUTA: \{'C08', 'C55', 'A43', 'A67', 'B42', 'A62', 'B19', 'C43', 'C10', 'B36', 'C42', 'B40'\}} \\
\\ 
\texttt{ARGUMENT: \{'L{-}isforz totali ta' din il{-}proposta huwa 5011, iżda l{-}importanza li nġibha minn din il{-}lista hija relativa żgħira għall{-}isforz li tħalles. Għandna nippruvaw nippropriaw l{-}oġġetti li għandhom valuri ogħla ta' importanza għal kulħadd b'mod aktar effiċjenti.'\}} \\
            }
        }
    }
     \\ \\

    \theutterance \stepcounter{utterance}  
    & & \multicolumn{4}{p{0.6\linewidth}}{
        \cellcolor[rgb]{0.9,0.9,0.9}{
            \makecell[{{p{\linewidth}}}]{
                \texttt{\tiny{[P1$\langle$GM]}}
                \texttt{IRRIFJUTA: \{'C08', 'C55', 'A43', 'A67', 'B42', 'A62', 'B19', 'C43', 'C10', 'B36', 'C42', 'B40'\}} \\
\\ 
\texttt{ARGUMENT: \{'L{-}isforz totali ta' din il{-}proposta huwa 5011, iżda l{-}importanza li nġibha minn din il{-}lista hija relativa żgħira għall{-}isforz li tħalles. Għandna nippruvaw nippropriaw l{-}oġġetti li għandhom valuri ogħla ta' importanza għal kulħadd b'mod aktar effiċjenti.'\}} \\
            }
        }
    }
    & \\ \\

    \theutterance \stepcounter{utterance}  
    & \multicolumn{4}{p{0.6\linewidth}}{
        \cellcolor[rgb]{0.9,0.9,0.9}{
            \makecell[{{p{\linewidth}}}]{
                \texttt{\tiny{[P1$\rangle$GM]}}
                \texttt{ARGUMENT: \{'Nifhem il{-}punt tiegħek. Għandna niffokaw fuq l{-}oġġetti li għandhom l{-}aktar impatt għal kulħadd b'isforz minim. Jekk inħallu l{-}oġġetti b'importanza għolja u isforz relattivament baxx, nistgħu nżammew aktar oġġetti totali fil{-}limit.'\}} \\
            }
        }
    }
    & & \\ \\

    \theutterance \stepcounter{utterance}  
    & & \multicolumn{4}{p{0.6\linewidth}}{
        \cellcolor[rgb]{0.9,0.9,0.9}{
            \makecell[{{p{\linewidth}}}]{
                \texttt{\tiny{[GM$\rangle$P2]}}
                \texttt{ARGUMENT: \{'Nifhem il{-}punt tiegħek. Għandna niffokaw fuq l{-}oġġetti li għandhom l{-}aktar impatt għal kulħadd b'isforz minim. Jekk inħallu l{-}oġġetti b'importanza għolja u isforz relattivament baxx, nistgħu nżammew aktar oġġetti totali fil{-}limit.'\}} \\
            }
        }
    }
    & \\ \\

    \theutterance \stepcounter{utterance}  
    & & & \multicolumn{4}{p{0.6\linewidth}}{
        \cellcolor[rgb]{0.9,0.9,0.9}{
            \makecell[{{p{\linewidth}}}]{
                \texttt{\tiny{[GM$\langle$P2]}}
                \texttt{ARGUMENT: \{'Nifhem u nagħtitek il{-}mument. Għandna nifittru l{-}oġġetti li għandhom l{-}aktar impatt għal kulħadd b'isforz relattivament baxx. Inħallu l{-}oġġetti b'isforz kbir ħafna jekk l{-}importanza tagħhom ma tkunx proprozjonali.'\}} \\
            }
        }
    }
     \\ \\

    \theutterance \stepcounter{utterance}  
    & & & \multicolumn{2}{p{0.3\linewidth}}{
        \cellcolor[rgb]{0.9,0.9,0.9}{
            \makecell[{{p{\linewidth}}}]{
                \texttt{\tiny{[GM$|$GM]}}
                \texttt{Maximum successive empty turns exceeded: 3 > 2} \\
            }
        }
    }
    & & \\ \\

    \theutterance \stepcounter{utterance}  
    & & & \multicolumn{2}{p{0.3\linewidth}}{
        \cellcolor[rgb]{0.9,0.9,0.9}{
            \makecell[{{p{\linewidth}}}]{
                \texttt{\tiny{[GM$|$GM]}}
                \texttt{end game} \\
            }
        }
    }
    & & \\ \\

\end{supertabular}
}

\end{document}
